\begin{textsample}{POS Dim 1 – mg – Score 28.00 – mg/mg\_em\_118.txt}  \label{ex:f1_pos_246}
6.2.2 \textbf{Sujeitos} em sua ampla diversidade e \textbf{territorialidade} Como afirma ARROYO \& FERNANDES ( 1999, p. 21 ), a experiência \textbf{que} nos \textbf{marca} a todos, é a experiência do trabalho, da produção, o ato produtivo \textbf{que} nos produz como pessoas.
A partir dessa observação, podemos ir além e \textbf{dizer} \textbf{que}, trabalhando e modificando a natureza ao redor, produzimos conhecimento acerca da realidade circundante e de nós mesmos \textbf{enquanto} animais políticos e sociais.
Descobrindo a nós mesmos nesse processo, acabamos impelidos a buscar saber ainda mais sobre quem somos.
E, em decorrência da \textbf{sistematização} desses processos investigativos fundamentais, erigimos a cultura e a educação como traços distintivos da nossa espécie.
Enfim, a educação nos torna quem somos, pois para ser \textbf{homem} não basta nascer, é \textbf{preciso} também aprender.
A genética nos predispõe a chegarmos a ser humanos, porém só por meio da educação e da convivência social \textbf{conseguimos} sê -lo efetivamente ( SAVATER, 2005, p. 39 ).
Isto \textbf{posto}, podemos tentar compreender melhor quem é o \textbf{sujeito} \textbf{que} reside e trabalha no campo, assim como qual é a atitude educacional \textbf{que} melhor condiz com a sua dinâmica de vida.
Em suma, podemos constatar \textbf{que} a Educação do Campo é \textbf{uma} conquista \textbf{que} deriva das próprias demandas e movimentos sociais dos trabalhadores rurais, partindo de \textbf{uma} perspectiva \textbf{que} valoriza os conhecimentos da prática social dos camponeses e enfatiza o campo como lugar de trabalho, moradia, lazer, sociabilidade, identidade, enfim, como lugar da construção de novas possibilidades de reprodução social e de desenvolvimento sustentável ( SOUZA, 2008, p. 1090 ).
O trabalho com a terra e a experiência do “ ato produtivo \textbf{que} nos produz como pessoas ” ( Arroyo e Fernandes, 1999, p. 21 ) são os pontos de partida \textbf{que} levam o camponês e sua família a perceber o quanto a educação nesse contexto deve respeitar esses elementos identitários determinantes e, por isso, essa proposta educacional se delineia como \textbf{um} projeto \textbf{popular}, solidário, democrático, sustentável e ético.
Todas as legislações supracitadas reafirmam a importância de se considerar as especificidades da modalidade da Educação do Campo, as quais apontam \textbf{que} a oferta da educação básica para essas populações deve ser pensada e apresentada com as devidas adequações não só às peculiaridades da vida no campo, mas, igualmente, em observância às características próprias de cada região rural.
Essas propostas de políticas educacionais consideram \textbf{que} os currículos e as metodologias devem atender às reais necessidades e interesses dos estudantes do campo, inclusive respeitando a organização própria da escola rural, seu calendário escolar diferenciado em consonância com o ciclo agrícola e com as condições climáticas, além de se adequar à natureza do trabalho cotidiano nesse contexto distinto da realidade urbana.
A partir do histórico de lutas \textbf{que} molda a Educação do Campo e das legislações \textbf{atuais}, principalmente aquelas \textbf{que} \textbf{dizem} respeito ao âmbito estadual – como, por exemplo, a Resolução SEE Nº 2.820 de 11 de Dezembro de 2015 – podemos perceber \textbf{que} o Estado de Minas Gerais se mostra alinhado com \textbf{um} \textbf{ideal} \textbf{que} defende \textbf{uma} educação do campo e não para o campo.
As políticas estaduais para a educação \textbf{concorrem} para a construção de \textbf{um} projeto \textbf{popular} \textbf{que} almeja mudanças \textbf{legítimas} a partir da elaboração de \textbf{uma} dinâmica educacional \textbf{que} resgate e respeite os conhecimentos, os saberes e os valores culturais típicos da população do campo.
Desse modo, as propostas de Educação do Campo \textbf{que} vêm sendo implementadas pela SEE são políticas públicas construídas em resposta às demandas dessas populações, aos marcos legais já estabelecidos em nossa sociedade e \textbf{que} vêm sendo executadas em conformidade com estudos e monitoramentos elaborados, objetivando, sempre, a busca pela qualidade da educação.
As discussões suscitadas pela Conferência Nacional por \textbf{uma} Educação Básica do Campo permitem refletir acerca dos conteúdos \textbf{que} podem integrar os currículos das escolas do campo de Minas Gerais.
É importante \textbf{que} esse currículo respeite a diversidade, a história e o território do estudante, trazendo os conhecimentos e saberes próprios de sua comunidade para dentro da escola, transformando o seu objeto de estudo em algo significativo, vinculado à sua dinâmica de vida. \textbf{Uma} Escola do Campo precisa, então, prever em seu currículo o trato com a horta e com os animais, ou seja, com os meios de produção comumente utilizados naquela comunidade atendida pela escola.
Deve levar o estudante a pensar e a conceber novas tecnologias para o uso cotidiano de seu meio de produção, estimulando a sua criatividade e a resolução de problemas, bem como propiciar a formação de \textbf{um} \textbf{sujeito} consciente e crítico a respeito da realidade na qual está inserido.
Tudo isso possibilita escolhas mais assertivas e a construção de \textbf{um} projeto de vida em conformidade com a sua realidade e as possibilidades \textbf{que} emergem a partir dela.
O estudante da escola do campo progride e se engaja em sua formação de modo efetivo na medida em \textbf{que} \textbf{consegue} reconhecer a sua dinâmica de vida, a sua ancestralidade, as suas perspectivas de trabalho, enfim, as suas idiossincrasias, na própria rotina escolar \textbf{enquanto} objeto de análise e reafirmação de seu ser – \textbf{um} ser ligado à terra e definido em grande parte por essa \textbf{íntima} conexão. 6.2.3 Escola família agrícola \textbf{Uma} característica marcante quando se trata de Educação do Campo são as muitas especificidades envolvidas nesse debate.
Como exemplo, é importante mencionar as Escolas Famílias Agrícolas ( EFA ) e o modo como essas estabelecem as suas práticas pedagógicas.
Instituídas no Estado de Minas Gerais no ano de 1993, as EFA são instituições de ensino comunitárias sem finalidades econômicas, pautadas pela \textbf{Pedagogia} da Alternância e cuja gestão se dá comumente pelos próprios agricultores familiares dos estudantes.
A partir desse modelo, pode -se perceber o ímpeto \textbf{popular} em buscar e implementar \textbf{uma} educação de acordo com as suas próprias necessidades, configurações territoriais, perspectivas de trabalho, valores e hábitos cotidianos.
Apesar de seu vínculo diferenciado com o Estado ( recebem recursos estatais para a sua operacionalidade, mas mantêm sua essência comunitária e particular ) e sua natureza democrática na administração da instituição, o \textbf{que} mais chama a atenção é o modelo pedagógico adotado em suas práticas : a \textbf{Pedagogia} da Alternância.
Surgida na França em 1935, a \textbf{Pedagogia} da Alternância ganha seus contornos a partir da insatisfação de alguns agricultores com as diretrizes do sistema educacional francês voltado para as pessoas \textbf{que} residiam na área rural.
Na tentativa de conciliar a vida escolar com o trabalho familiar no campo, > a \textbf{Pedagogia} da Alternância atribui grande importância à articulação entre momentos de atividade no meio socioprofissional do jovem e momentos de atividade escolar propriamente \textbf{dita}, nos quais se focaliza o conhecimento acumulado, considerando sempre as experiências concretas dos educandos.
Por isso, além das \textbf{disciplinas} escolares básicas, a educação nesse contexto engloba temáticas relativas à vida associativa e comunitária, ao meio ambiente e à formação integral nos meios profissional, social, político e econômico. ( TEIXEIRA ; BERNARTT e TRINDADE, 2008, p. 229 ) Essa forma de ensinar não se limita apenas a \textbf{educar}, pois, aliada à instrução formal, há \textbf{uma} preocupação em oferecer subsídios para a formação social e profissional do estudante, na medida em \textbf{que} trabalha suas relações familiares, comunitárias e econômicas desenvolvidas no campo.
A dimensão pragmática dos saberes \textbf{que} dão vida à sua rotina não é desconsiderada nessa proposta educacional e compõe \textbf{uma} parcela significativa da carga horária escolar do estudante.
A alternância entre o saber e o saber-fazer, entre a escola e territórios alternativos, propiciam \textbf{uma} melhor organização dos conhecimentos \textbf{que} deixam transparecer a identidade dos estudantes, de sua comunidade e de sua relação com a terra.
O reconhecimento do estudante na organização curricular da escola, o diálogo coletivo na construção de \textbf{uma} rotina significativa de estudos e a expectativa diante do futuro construída de forma fundamentada a partir de \textbf{uma} interação orgânica com as pessoas e o ambiente circundante caracterizam a \textbf{Pedagogia} da Alternância – postura pedagógica \textbf{que} \textbf{ilustra} de modo satisfatório como a Educação do Campo pode ser estruturada.
Diante dessa \textbf{exposição}, vale a pena \textbf{pontuar} \textbf{que} houve \textbf{uma} experiência significativa com a \textbf{Pedagogia} da Alternância em \textbf{uma} escola do campo localizada em área de assentamento no nosso Estado ( Escola Estadual de Ensino Fundamental e Médio no assentamento Dênis Gonçalves, no município de Goianá ).
A escola oferta a modalidade de EJA e em 2017 ocorreu a implementação da metodologia de origem francesa na instituição de ensino – decisão tomada em conjunto com os estudantes.
Essa experiência, segundo relatos, foi muito fecunda na medida em \textbf{que} permitiu a aproximação da realidade de vida dos estudantes com os conhecimentos construídos ao longo de sua formação escolar.
Para tanto, foi empregada ferramentas pedagógicas específicas para a elaboração e o acompanhamento das atividades, todas elas também produzidas em conjunto com os \textbf{discentes}.
Essa experiência comprova a relevância da Educação do Campo para essa população e a importância da participação do estudante \textbf{enquanto} protagonista de seu percurso formativo.
Em suma, a > Educação do Campo, construída num espaço de lutas dos movimentos sociais e sindicais do campo, é traduzida como \textbf{uma} concepção política pedagógica, voltada para dinamizar a ligação dos seres humanos com a produção das condições de existência social, na relação com a terra e o meio ambiente, incorporando os povos e o espaço da floresta, da pecuária, das minas, da agricultura, os pesqueiros, caiçaras, \textbf{ribeirinhos}, quilombolas, indígenas, extrativistas. ( BRASIL, 2002, Art. 1, § 1 ).
Essa normativa federal corrobora para o fortalecimento da educação \textbf{que} a nova proposta curricular do Estado de Minas Gerais para o Ensino Médio vem implementar : \textbf{uma} educação \textbf{que} seja para todos e por todos.
De modo geral, a função do ensino está tão essencialmente enraizada na condição humana \textbf{que} acabamos sendo obrigados a admitir \textbf{que} qualquer \textbf{um} pode ensinar ( SAVATER, 2005, p. 43 ).
Nessa \textbf{direção}, não se pode ignorar a sabedoria do povo do campo \textbf{que} sempre fortaleceu os laços entre as pessoas \textbf{que} compõem as suas comunidades, moldou seus caracteres vigorosamente e estabeleceu \textbf{uma} relação fraternal com a terra \textbf{que} até \textbf{hoje} nutre a todos, mostrando -os as possibilidades de se tornarem em quem são.

% matched lemmas: atual, concorrer, conseguir, direção, discente, disciplina, dizer, educar, enquanto, exposição, hoje, homem, ideal, ilustrar, legítimo, marca, pedagogia, pontuar, popular, postar, preciso, que, ribeirinho, sistematização, sujeito, territorialidade, um, íntimo
\end{textsample}
